% =============================================================================
% CHECKLIST DI RIPRODUCIBILITA FACOLTATIVA - ITALIANO
% Adattamento per tesi del formato della NeurIPS Paper Checklist:
% Si / No / N/A con evidenza o breve giustificazione.
% Formato di riferimento: https://neurips.cc/public/guides/PaperChecklist
% Non e il modulo ufficiale NeurIPS e non implica conformita alla conferenza.
% Disabilitarla in configuration.tex con:
% \providecommand{\ISISReproducibilityChecklistChoice}{disabled}
% =============================================================================

\chapter{Checklist di riproducibilit\`a}\label{app:reproducibility-checklist}

\begin{keypoint}[Facoltativa e basata sull'applicabilit\`a]
Questa checklist \`e uno strumento facoltativo per la tesi. Conservare solo le
voci pertinenti e indicare \textbf{N/A} quando una voce non \`e applicabile o
l'informazione richiesta non \`e disponibile. Per ogni risposta, citare una
sezione della tesi, un'appendice, un file del repository, una scheda del dataset
oppure fornire una breve giustificazione. Il formato \`e adattato dalla checklist
NeurIPS e non costituisce un modulo ufficiale di sottomissione a una conferenza.
\end{keypoint}

\noindent\textbf{Uso suggerito.}
Compilare la checklist dopo la bibliografia e l'eventuale appendice tecnica.
Rivederla con il relatore prima della consegna e rimuovere collegamenti o
informazioni riservate che non devono comparire nella versione pubblica.

\begin{reproducibilitychecklist}

\reproitem{Affermazioni e ambito}
{Il sommario e l'introduzione descrivono correttamente contributo, ambito,
ipotesi e livello di generalit\`a supportato dai risultati?}
{Citare il sommario, l'elenco dei contributi e le sezioni contenenti le evidenze principali.}

\reproitem{Limiti}
{Sono discussi i limiti importanti, le ipotesi forti, i possibili casi di
fallimento e i confini di validit\`a?}
{Citare una sezione sui limiti o sulle minacce alla validit\`a e motivare eventuali omissioni.}

\reproitem{Teoria, ipotesi e dimostrazioni}
{Quando sono presenti risultati teorici, tutte le ipotesi sono dichiarate e le
dimostrazioni complete o i relativi riferimenti sono forniti?}
{Indicare N/A per lavori senza affermazioni teoriche; altrimenti citare teoremi e appendici.}

\reproitem{Percorso di riproduzione}
{Esiste un percorso pratico che consenta a un altro ricercatore di riprodurre o
verificare in modo indipendente il contributo principale?}
{Indicare appendice, repository, accesso al modello, protocollo o altro meccanismo di verifica.}

\reproitem{Accesso a codice e dati}
{Quando sono stati eseguiti esperimenti, codice, dati, artefatti del modello e
istruzioni di esecuzione sono disponibili, oppure le restrizioni sono spiegate?}
{Fornire link persistenti o percorsi, versioni, condizioni di accesso e motivo delle restrizioni.}

\reproitem{Configurazione sperimentale}
{Sono specificati dataset, split, preprocessing, baseline, iperparametri,
procedure di selezione, seed casuali, criteri di arresto e versioni software?}
{Citare metodologia, tabella in appendice, file di configurazione e ambiente o lock file.}

\reproitem{Incertezza ed evidenza statistica}
{Sono riportate esecuzioni ripetute, barre di errore, intervalli di confidenza,
test statistici o altre misure appropriate di incertezza, definite correttamente?}
{Indicare cosa varia tra le esecuzioni, il numero di ripetizioni e il metodo di calcolo.}

\reproitem{Risorse di calcolo}
{Sono documentati tipo di hardware, memoria, tempo di esecuzione, storage e
stima del calcolo totale necessario per gli esperimenti riportati?}
{Citare la descrizione dell'ambiente e distinguere, se rilevante, esperimenti finali ed esplorativi.}

\reproitem{Etica della ricerca}
{Il lavoro \`e stato verificato rispetto ai requisiti applicabili di integrit\`a
scientifica, sicurezza, privacy ed etica professionale?}
{Citare dichiarazione, procedura di revisione o regola istituzionale, oppure motivare la non applicabilit\`a.}

\reproitem{Impatto sociale}
{Quando pertinente, sono discussi usi dannosi, equit\`a, privacy, sicurezza e
possibili conseguenze del deployment?}
{Citare la discussione di impatto o rischio e le eventuali misure di mitigazione.}

\reproitem{Misure di salvaguardia}
{Per artefatti con rischio significativo di abuso o duplice uso, sono descritte
misure proporzionate di rilascio, condizioni d'uso o controlli di accesso?}
{Indicare N/A se non esiste un rischio elevato; altrimenti descrivere salvaguardie e rischio residuo.}

\reproitem{Licenze e condizioni d'uso}
{Dataset, codice, modelli e altri asset di terzi sono citati, versionati e usati
nel rispetto delle licenze e delle condizioni applicabili?}
{Elencare licenza o condizioni per ogni asset importante e segnalare incertezze irrisolte.}

\reproitem{Documentazione dei nuovi artefatti}
{Se la tesi rilascia nuovo codice, dati, modelli o benchmark, sono documentati
scopo, contenuto, provenienza, limiti, licenza e informazioni di manutenzione?}
{Indicare N/A se non vengono rilasciati artefatti; altrimenti citare README, data card o model card.}

\reproitem{Partecipanti umani o crowdsourcing}
{Se persone hanno fornito dati o partecipato allo studio, reclutamento,
istruzioni, consenso, compenso e procedure sono documentati?}
{Indicare N/A se non sono stati coinvolti partecipanti umani o crowd worker.}

\reproitem{Approvazione etica o istituzionale}
{Quando richiesto, i rischi per i partecipanti sono stati valutati ed \`e stata
ottenuta l'approvazione o l'esenzione dall'organo competente?}
{Fornire il riferimento non riservato dell'approvazione o motivare perch\'e non era necessaria.}

\reproitem{Uso di IA generativa o LLM}
{Quando un sistema generativo o un LLM costituisce una parte importante,
originale o non standard del metodo, ruolo, versione, configurazione e verifica
sono descritti?}
{Indicare N/A per semplice assistenza di scrittura o formattazione che non modifica metodo o evidenze.}

\end{reproducibilitychecklist}

\begin{keypoint}[Revisione finale]
Una risposta ``No'' non \`e automaticamente un difetto: deve essere accompagnata
da una spiegazione chiara e, quando possibile, da una mitigazione o da un'azione
futura. Rimuovere le risposte segnaposto e verificare che ogni sezione, file e
collegamento citato esista nella versione consegnata.
\end{keypoint}
